% ========================================================
% BÀI TẬP TOÁN 5 - KẾT NỐI TRI THỨC VỚI CUỘC SỐNG
% Bài 38: Tìm hai số khi biết tổng và tỉ số
% Định dạng: BTTeX Part III - Trắc nghiệm điền số ngắn
% ========================================================

% Câu 1
\begin{ex}
  Hai số có tổng là 45 và tỉ số là $\frac{1}{2}$. Tìm số bé.
  \shortans{15}
  \loigiai{
    Tổng số phần bằng nhau là 1 + 2 = 3 phần. Giá trị một phần là 45 : 3 = 15. Số bé là 15.
  }
\end{ex}

% Câu 2
\begin{ex}
  Hai số có tổng là 45 và tỉ số là $\frac{1}{2}$. Tìm số lớn.
  \shortans{30}
  \loigiai{
    Tổng số phần bằng nhau là 1 + 2 = 3 phần. Giá trị một phần là 45 : 3 = 15. Số lớn là 30.
  }
\end{ex}

% Câu 3
\begin{ex}
  Hai số có tổng là 72 và tỉ số là $\frac{3}{5}$. Tìm số bé.
  \shortans{27}
  \loigiai{
    Tổng số phần bằng nhau là 3 + 5 = 8 phần. Giá trị một phần là 72 : 8 = 9. Số bé là 27.
  }
\end{ex}

% Câu 4
\begin{ex}
  Hai số có tổng là 72 và tỉ số là $\frac{3}{5}$. Tìm số lớn.
  \shortans{45}
  \loigiai{
    Tổng số phần bằng nhau là 3 + 5 = 8 phần. Giá trị một phần là 72 : 8 = 9. Số lớn là 45.
  }
\end{ex}

% Câu 5
\begin{ex}
  Hai số có tổng là 100 và tỉ số là $\frac{1}{3}$. Tìm số bé.
  \shortans{25}
  \loigiai{
    Tổng số phần bằng nhau là 1 + 3 = 4 phần. Giá trị một phần là 100 : 4 = 25. Số bé là 25.
  }
\end{ex}

% Câu 6
\begin{ex}
  Hai số có tổng là 100 và tỉ số là $\frac{1}{3}$. Tìm số lớn.
  \shortans{75}
  \loigiai{
    Tổng số phần bằng nhau là 1 + 3 = 4 phần. Giá trị một phần là 100 : 4 = 25. Số lớn là 75.
  }
\end{ex}

% Câu 7
\begin{ex}
  Hai số có tổng là 120 và tỉ số là $\frac{2}{3}$. Tìm số bé.
  \shortans{48}
  \loigiai{
    Tổng số phần bằng nhau là 2 + 3 = 5 phần. Giá trị một phần là 120 : 5 = 24. Số bé là 48.
  }
\end{ex}

% Câu 8
\begin{ex}
  Hai số có tổng là 120 và tỉ số là $\frac{2}{3}$. Tìm số lớn.
  \shortans{72}
  \loigiai{
    Tổng số phần bằng nhau là 2 + 3 = 5 phần. Giá trị một phần là 120 : 5 = 24. Số lớn là 72.
  }
\end{ex}

% Câu 9
\begin{ex}
  Hai số có tổng là 150 và tỉ số là $\frac{1}{4}$. Tìm số bé.
  \shortans{30}
  \loigiai{
    Tổng số phần bằng nhau là 1 + 4 = 5 phần. Giá trị một phần là 150 : 5 = 30. Số bé là 30.
  }
\end{ex}

% Câu 10
\begin{ex}
  Hai số có tổng là 150 và tỉ số là $\frac{1}{4}$. Tìm số lớn.
  \shortans{120}
  \loigiai{
    Tổng số phần bằng nhau là 1 + 4 = 5 phần. Giá trị một phần là 150 : 5 = 30. Số lớn là 120.
  }
\end{ex}

% Câu 11
\begin{ex}
  Hai số có tổng là 84 và tỉ số là $\frac{3}{4}$. Tìm số bé.
  \shortans{36}
  \loigiai{
    Tổng số phần bằng nhau là 3 + 4 = 7 phần. Giá trị một phần là 84 : 7 = 12. Số bé là 36.
  }
\end{ex}

% Câu 12
\begin{ex}
  Hai số có tổng là 84 và tỉ số là $\frac{3}{4}$. Tìm số lớn.
  \shortans{48}
  \loigiai{
    Tổng số phần bằng nhau là 3 + 4 = 7 phần. Giá trị một phần là 84 : 7 = 12. Số lớn là 48.
  }
\end{ex}

% Câu 13
\begin{ex}
  Hai số có tổng là 96 và tỉ số là $\frac{1}{5}$. Tìm số bé.
  \shortans{16}
  \loigiai{
    Tổng số phần bằng nhau là 1 + 5 = 6 phần. Giá trị một phần là 96 : 6 = 16. Số bé là 16.
  }
\end{ex}

% Câu 14
\begin{ex}
  Hai số có tổng là 96 và tỉ số là $\frac{1}{5}$. Tìm số lớn.
  \shortans{80}
  \loigiai{
    Tổng số phần bằng nhau là 1 + 5 = 6 phần. Giá trị một phần là 96 : 6 = 16. Số lớn là 80.
  }
\end{ex}

% Câu 15
\begin{ex}
  Hai số có tổng là 110 và tỉ số là $\frac{4}{7}$. Tìm số bé.
  \shortans{40}
  \loigiai{
    Tổng số phần bằng nhau là 4 + 7 = 11 phần. Giá trị một phần là 110 : 11 = 10. Số bé là 40.
  }
\end{ex}

% Câu 16
\begin{ex}
  Hai số có tổng là 110 và tỉ số là $\frac{4}{7}$. Tìm số lớn.
  \shortans{70}
  \loigiai{
    Tổng số phần bằng nhau là 4 + 7 = 11 phần. Giá trị một phần là 110 : 11 = 10. Số lớn là 70.
  }
\end{ex}

% Câu 17
\begin{ex}
  Hai số có tổng là 132 và tỉ số là $\frac{5}{6}$. Tìm số bé.
  \shortans{60}
  \loigiai{
    Tổng số phần bằng nhau là 5 + 6 = 11 phần. Giá trị một phần là 132 : 11 = 12. Số bé là 60.
  }
\end{ex}

% Câu 18
\begin{ex}
  Hai số có tổng là 132 và tỉ số là $\frac{5}{6}$. Tìm số lớn.
  \shortans{72}
  \loigiai{
    Tổng số phần bằng nhau là 5 + 6 = 11 phần. Giá trị một phần là 132 : 11 = 12. Số lớn là 72.
  }
\end{ex}

% Câu 19
\begin{ex}
  Hai số có tổng là 64 và tỉ số là $\frac{1}{7}$. Tìm số bé.
  \shortans{8}
  \loigiai{
    Tổng số phần bằng nhau là 1 + 7 = 8 phần. Giá trị một phần là 64 : 8 = 8. Số bé là 8.
  }
\end{ex}

% Câu 20
\begin{ex}
  Hai số có tổng là 64 và tỉ số là $\frac{1}{7}$. Tìm số lớn.
  \shortans{56}
  \loigiai{
    Tổng số phần bằng nhau là 1 + 7 = 8 phần. Giá trị một phần là 64 : 8 = 8. Số lớn là 56.
  }
\end{ex}

% Câu 21
\begin{ex}
  Hai số có tổng là 144 và tỉ số là $\frac{1}{2}$. Tìm số bé.
  \shortans{48}
  \loigiai{
    Tổng số phần bằng nhau là 1 + 2 = 3 phần. Giá trị một phần là 144 : 3 = 48. Số bé là 48.
  }
\end{ex}

% Câu 22
\begin{ex}
  Hai số có tổng là 144 và tỉ số là $\frac{1}{2}$. Tìm số lớn.
  \shortans{96}
  \loigiai{
    Tổng số phần bằng nhau là 1 + 2 = 3 phần. Giá trị một phần là 144 : 3 = 48. Số lớn là 96.
  }
\end{ex}

% Câu 23
\begin{ex}
  Hai số có tổng là 200 và tỉ số là $\frac{3}{7}$. Tìm số bé.
  \shortans{60}
  \loigiai{
    Tổng số phần bằng nhau là 3 + 7 = 10 phần. Giá trị một phần là 200 : 10 = 20. Số bé là 60.
  }
\end{ex}

% Câu 24
\begin{ex}
  Hai số có tổng là 200 và tỉ số là $\frac{3}{7}$. Tìm số lớn.
  \shortans{140}
  \loigiai{
    Tổng số phần bằng nhau là 3 + 7 = 10 phần. Giá trị một phần là 200 : 10 = 20. Số lớn là 140.
  }
\end{ex}

% Câu 25
\begin{ex}
  Hai số có tổng là 90 và tỉ số là $\frac{2}{7}$. Tìm số bé.
  \shortans{20}
  \loigiai{
    Tổng số phần bằng nhau là 2 + 7 = 9 phần. Giá trị một phần là 90 : 9 = 10. Số bé là 20.
  }
\end{ex}
